\chapter{Related work and positioning}
\label{chap:related}

This chapter situates our work in the deep-learning literature for time series,
limited to the ideas we actually build on, and states our positioning at the
end.

\section{Deep learning for time-series classification}
\label{sec:rw-tsc}

Time-series classification (TSC) has shifted from distance-based methods
(e.g.\ nearest neighbour with dynamic time warping) to deep learning. The review
of Ismail Fawaz et al.\ \cite{intfawaz2019deep} benchmarks the main families and
identifies fully-convolutional and residual networks as the strongest
general-purpose baselines, a conclusion already drawn by Wang et al.\
\cite{wang2017time}, whose end-to-end fully-convolutional network (FCN) and 1D
ResNet outperform most hand-crafted pipelines without feature engineering. These
two works directly motivate our choice of a convolutional baseline and a
residual variant. Temporal convolutional networks \cite{bai2018empirical} show,
more generally, that convolutions with a sufficient receptive field match or beat
recurrent models on sequence tasks while being cheaper to train --- relevant to
an operational beacon with a limited compute budget.

\section{Convolutional and residual architectures}
\label{sec:rw-cnn}

The convolutional building block goes back to LeCun et al.\
\cite{lecun1998gradient}: weight-shared local filters give translation
equivariance and a parameter count independent of the input length, which is
exactly what is needed to detect an anomaly that may appear anywhere in a window.
Very deep stacks are hard to optimise because gradients vanish; He et al.\
\cite{he2016deep} solve this with \emph{residual} connections that let a block
default to the identity, making depth essentially free. Training is further
stabilised by normalisation: batch normalisation \cite{ioffe2015batch} reduces
internal covariate shift but couples train and inference statistics, whereas
group normalisation \cite{wu2018group} computes statistics per sample and so
behaves identically at train and inference time --- a property we exploit for a
signal whose amplitude scale drifts.

\section{Attention and Transformer encoders}
\label{sec:rw-attn}

Self-attention \cite{vaswani2017attention} lets every position in a sequence
attend to every other with a learnt, content-dependent weighting, giving an
unbounded receptive field in a single layer; layer normalisation
\cite{ba2016layer} and a learnable aggregation token \cite{devlin2019bert} make
such encoders practical for classification. Hybrids that place a Transformer
encoder on top of a convolutional stem are now common for noisy physiological
signals: Hu and Chen \cite{hu2022eeg}, cited in the project brief, combine
convolutions with multi-head self-attention for EEG emotion recognition. Our
third architecture follows the same recipe, transposed to gamma-dose signals.

\section{Anomaly detection and explainability}
\label{sec:rw-ad}

Our task is closer to \emph{classification of known morphologies} than to
unsupervised anomaly detection, the latter surveyed by Bl\'azquez-Garc\'ia et
al.\ \cite{blazquezgarcia2021review}; we borrow from it the event-level,
detection-oriented evaluation rather than a purely point-wise accuracy. Finally,
because an operational alert must be trustworthy, explainability matters: the
brief points to Grad-CAM \cite{selvaraju2017gradcam}, which produces
gradient-based saliency over the input. Our convolutional models expose the hooks
required for it (\cref{chap:conclusion}), but a full explainability study is out
of scope here.

\section{Positioning}
\label{sec:rw-position}

We adopt established architectures --- a convolutional baseline, a residual
network \cite{he2016deep,wang2017time}, and a CNN--Transformer hybrid
\cite{vaswani2017attention,hu2022eeg} --- rather than proposing a new one. The
specific contribution of this project is therefore not architectural but
methodological and applicative: (i) a \emph{simplified yet statistically
calibrated} synthetic generator for noisy gamma-dose signals with controllable
anomaly morphologies, and (ii) a controlled, like-for-like comparison of the
three architectures along the three axes named in the brief --- classification
performance, computational cost, and robustness to distribution shift. To our
knowledge this combination has not been reported for radiological-beacon anomaly
classification.
